\section{Further research questions and concrete next targets}
\label{ag:sec:questions}

The following problems continue the exact-identity programme rather
than asking for bounds or extrapolating a numerical pattern. The
proved formulas give initial conditions and explicit checks for each.

\paragraph{1. Complete the coupled next Gamma sector.}
Find a finite expression for the coefficient
\[
 [u_0^a u_1^k u_2^b u_3^2]N_{a+2k+3b+8}
\]
for every $b>0$.
The present theorem supplies $b=0$, while the predecessor supplies the
pure $u_2$ projection. A quadratic background for $u_0$ would retain
$u_2$ but replace constant-coefficient weighted adjoints by variable-
coefficient ones. The precise missing step is a canonical symbol
calculus compatible with that background and the original normal basis.
A formula should include the parity zeros and the retained endpoint
primitive, not just the first few nonzero weights.

\paragraph{2. Determine the first genuinely nonlinear polarization.}
The affine quadratic symbol is a polynomial in one invariant $\eta$.
At cubic and quartic perturbation order there are several symmetric
invariants after imposing $\sum\xi_i+\kappa=0$. Identify an invariant
basis which maps triangularly to the functional-monomial normal form.
Use it to compute complete families with three or four high-derivative
factors, distinguishing cubic terms such as $u_2^3$ from quartic terms
such as $u_2^2u_3^2$. The source's all-$u_2$ formulas provide stringent
independent checks but not the mixed-sector solution.

\paragraph{3. Classify elementary symmetric profiles.}
For which even analytic profiles $K$ do all square-sector kernels
$g_{K,r}$ belong to a fixed finite combination of hyperbolic functions
and their divided differences? The finite Gamma mixtures in
Section~\ref{ag:sec:profiles} give one class. Seek a converse under
explicit growth and analyticity hypotheses, or exhibit another class
coming from an interpretable Gamma or Barnes ratio. A claim about
minimality must specify the allowed coefficient field and operations.

\paragraph{4. Generalize the selector classification to other intervals.}
On $[1,3]$, integer translation leaves only the bases $1$ and $2^{-s}$
besides shifted Riemann zeta functions. For intervals $[A,B]$ with
positive integer endpoints, prescribe the finite exponential remainder
and characterize weights selecting exactly one shifted zeta function.
For rational endpoints the natural coordinates are Hurwitz values or
Dirichlet $L$-functions. The goal is an exact functional classification,
including endpoint multiplicities and a proof of the necessary
Dirichlet-series independence statement.

\paragraph{5. Find Lerch selectors with a prescribed polylogarithm order.}
The general Lerch endpoint formula is finite, but the rational factors
$z^{-A}$ prevent the simple odd-derivative diagonalization used at $z=1$.
Characterize polynomial weights, possibly depending rationally on $z$,
whose transform has one polylogarithm order plus an elementary
remainder. Determine whether the required endpoint conditions form a
finite interpolation problem, and prove its confluence as $z\to1$.
The explicit quartic identity in Section~\ref{ag:sec:polylogs} is a
normalization test rather than a general classification.

\paragraph{6. Put all selector Stieltjes orders in a uniform closed form.}
Equation~\eqref{ag:eq:selector-Stieltjes} is a finite coefficient rule
at every order. Seek a compact bivariate generating identity in the
selector degree and Stieltjes order that preserves the single spectral
center and makes the harmonic coefficients explicit. Determine which
anchored integrals in the shift can be expressed through a fixed family
of multiple-Gamma or zeta-derivative functions. Any such primitive must
retain the anchor conditions from Section~\ref{ag:sec:stieltjes}.

\paragraph{7. Describe the numerical image of the Gamma square sector.}
The coefficient theorem is formal. Evaluate or systematically reduce
the ordinary shifted integrals and the original-coordinate finite parts
of $\psi^a(\psi')^{2k}(\psi^{(r)})^2$ in a declared basis of convergent
multiple-zeta derivatives or polylogarithmic integrals. The derivative-
ideal pole calculus in the manuscript supplies convergent coordinates,
but it does not establish their arithmetic minimality. Retain the
primitive endpoint term when transporting the result to inverse-Gamma
moments.

\paragraph{8. Return to the unequal-tail boundary determinant.}
The all-$r$ Gamma theorem does not prove the unequal-even-tail
optimal-prefix conjecture. The next unresolved general fixed-gap
problem in the inspected predecessor is the three-column boundary
determinant for gap six. Find an exact scalar recurrence, a determinant
identity, or an arithmetic normalization that proves nonvanishing.
Entrywise positivity and pairwise ratio separation are not sufficient.
Any resulting inverse row can then be combined with the predecessor's
already proved universal finite telescoper.

\paragraph{9. Build a proof-carrying coefficient compiler.}
Formalize the terminating normal-form algorithm and its identity
$P=\NF(P)+\D\Prim(P)$ in a proof assistant. Add the Appell covariance
lemma and the weighted quadratic symbol as reusable algebraic results.
A useful initial target is a checker for the supplied rational JSON
certificates, followed by the all-order coefficient proof. Certified
real or complex interval checks of the displayed analytic integrals
would be a separate numerical project, not a replacement for analytic
continuation proofs.

\paragraph{10. Preserve the original period questions.}
The Gaussian $S_6$ and current revised $S_8$ candidates call for exact
polylogarithmic identities with fixed branches and complete lower-weight
constants. Neither selector closure nor normal-form coefficient
closure settles them. A promising next step is a parameter identity
whose rigorously controlled boundary value matches the precise frozen
candidate, rather than a new unconstrained integer-relation search.

\subsection*{Conclusion}
An affine perturbation converts an infinite family of Gamma normal-form
coefficients into one entire divided difference. The beta representation
of its coefficients supplies a common route to finite Hurwitz and Lerch
endpoint formulas, generalized Stieltjes averages, and anchored
primitives. Bernoulli endpoint data then diagonalize the integer-
interval arithmetic and give a genuine single-zeta functional
classification. These are exact identities with separate formal,
analytic, and computational meanings. Their limitations identify the
next nonlinear and arithmetic questions without weakening the proved
all-index statements.
